\section*{Capítulo \ref{ch:b2:blood-pressure} --- Regulação da pressão arterial e exercício}

\begin{solution}{exo:b2:blood-pressure:1}
$\bar P - P_{\text{ven}} = Q\,R = f\,V_s\,R$. Frequência cardíaca (o \omterm{prop:b2:heart:conduction}{nó
sinoatrial}, sob o vago e os nervos simpáticos); \omterm{prop:b2:heart:cycle}{volume sistólico} (o
\omterm{def:b2:heart:anatomy}{ventrículo}: seu enchimento, fixado pelas \omterm{def:b2:blood-circulation:vessels}{veias}, e sua contratilidade,
fixada pelos nervos simpáticos); resistência (o músculo liso
arteriolar, sob os nervos simpáticos, os hormônios e os metabólitos
locais).
\end{solution}

\begin{solution}{exo:b2:blood-pressure:2}
Sensores: receptores de estiramento nos seios carotídeos e no arco
aórtico. Centro: os centros cardiovasculares do bulbo. Efetores: o vago
(frequência) e os nervos simpáticos (frequência, contratilidade, tônus
arteriolar e venoso). Retroalimentação negativa, cerca de um segundo de
atraso. Sem os sensores, a pressão média ao longo de um dia permanece
normal, mas a pressão de minuto a minuto oscila violentamente a cada
mudança de postura e a cada emoção.
\end{solution}

\begin{solution}{exo:b2:blood-pressure:3}
O rim (células arteriolares) libera renina quando sua perfusão ou o
aporte de sódio cai ou quando seus nervos simpáticos disparam; a renina
corta o angiotensinogênio (fígado) em angiotensina I; uma enzima
pulmonar a converte em angiotensina II, que contrai as \omterm{def:b2:blood-circulation:vessels}{arteríolas} e
estimula o córtex suprarrenal a secretar aldosterona; a aldosterona faz
o rim reter sódio e água, elevando o volume sanguíneo.
\end{solution}

\begin{solution}{exo:b2:blood-pressure:4}
A dilatação metabólica das próprias \omterm{def:b2:blood-circulation:vessels}{arteríolas} do músculo por seus
metabólitos; um \omterm{thm:b2:heart:starling}{débito cardíaco} cinco vezes maior, impulsionado pelos
nervos simpáticos e pela bomba muscular; a constrição simpática das
\omterm{def:b2:blood-circulation:vessels}{arteríolas} do intestino, dos rins e dos músculos em repouso, que cedem
fluxo.
\end{solution}

\begin{solution}{exo:b2:blood-pressure:5}
$83\times 1.2 = \qty{100}{mmHg}$; $333\times 0.35 = \qty{117}{mmHg}$.
Resistência reduzida por um fator $1.2/0.35 = 3.4$: raio aumentado por
um fator $3.4^{1/4} =
1.36$.
\end{solution}

\begin{solution}{exo:b2:blood-pressure:6}
$30/(1 + G)$: 10, 6 e \qty{3.3}{mmHg}. Com metade do ganho, a queda
residual dobra: ela fica tonta ao se levantar, vê tudo cinza e pode
desmaiar.
\end{solution}

\begin{solution}{exo:b2:blood-pressure:7}
$\rho g h$: $1060\times 9.81\times 0.4 = \qty{4160}{Pa} =
\qty{31}{mmHg}$; $1060\times 9.81\times 1.2 = \qty{12500}{Pa} =
\qty{94}{mmHg}$. Encéfalo $95 - 31 = \qty{64}{mmHg}$; tornozelo
$95 + 94 =
\qty{189}{mmHg}$. A cabeça de uma girafa a \qty{2.5}{m} custa
\qty{195}{mmHg}: ela precisa de uma pressão média de uns
\qty{250}{mmHg} no \omterm{def:b2:heart:anatomy}{coração}, e de paredes vasculares espessas e de uma
pele justa nas pernas.
\end{solution}

\begin{solution}{exo:b2:blood-pressure:8}
$Q = 280/(190 - 130) = \qty{4.7}{L/min}$. Atleta: $5000/(200 - 30) =
\qty{29}{L/min}$; \omterm{prop:b2:heart:cycle}{volume sistólico} $\num{29400}/190 = \qty{155}{mL}$.
\end{solution}

\begin{solution}{exo:b2:blood-pressure:9}
$35/5 = \qty{7}{mmHg}$ a menos: cerca de \qty{88}{mmHg}. Restauração do
volume: reabsorção de líquido intersticial para os \omterm{def:b2:blood-circulation:vessels}{capilares} (de
minutos a uma hora); redução da urina sob o \omterm{prop:b2:blood-pressure:raas}{hormônio antidiurético} e a
aldosterona (de horas a dias); sede e ingestão de água (horas); \omterm{def:b2:blood-circulation:blood}{hemácias}
da medula óssea (semanas). A pele fica pálida e fria porque o reflexo
fechou suas \omterm{def:b2:blood-circulation:vessels}{arteríolas} para reservar a pressão ao encéfalo e ao
\omterm{def:b2:heart:anatomy}{coração}.
\end{solution}

\begin{solution}{exo:b2:blood-pressure:10}
Músculo com $R = 0.5$: $95/0.5 = \qty{190}{mL/s}$, \qty{11.4}{L/min};
com os \qty{4}{L/min} dos outros órgãos, o \omterm{def:b2:heart:anatomy}{coração} precisaria de
\qty{15.4}{L/min}. Com $Q$ fixo em \qty{5}{L/min}, a resistência total
cai de 1.14 para $1/(1/1.14 - 1/5 + 1/0.5) =
\qty{0.37}{mmHg.s/mL}$ e a pressão para \qty{31}{mmHg} --- o encéfalo
desmaiaria. O corpo eleva o débito na direção do primeiro caso e
contrai os outros leitos vasculares, de modo que a pressão se mantém e
o músculo recebe seu fluxo.
\end{solution}

\begin{solution}{exo:b2:blood-pressure:11}
(a) O rim mal irrigado leu uma pressão baixa, liberou renina, e a
angiotensina e a aldosterona elevaram a resistência e o volume. (b) O
\omterm{def:b2:blood-pressure:baroreflex}{barorreflexo} amortece as mudanças, mas se reajusta a qualquer nível que
persista, e não consegue alterar a retenção de volume pelo rim. (c) A
fonte de renina, e o órgão que exigia uma pressão mais alta, tinham
desaparecido. (d) Teria bloqueado a conversão em angiotensina II:
resistência menor, menos aldosterona, pressão mais baixa.
\end{solution}

\begin{solution}{exo:b2:blood-pressure:12}
O \omterm{def:b2:blood-pressure:baroreflex}{barorreflexo} age em um segundo, com um ganho de cerca de quatro e um
valor de referência que deriva para a pressão predominante: ele
amortece as flutuações e não consegue manter um nível. O rim age ao
longo de horas e dias, com um ganho na prática infinito --- continua
excretando ou retendo até atingir a pressão para a qual está ajustado
---, e seu valor de referência é fixado por sua própria fisiologia: ele
mantém o nível. Uma elevação crônica significa, portanto, que o valor
de referência do rim se deslocou, e só os fármacos que agem sobre a alça
do rim ou sobre a resistência que ele exige a fazem baixar.
\end{solution}

\begin{solution}{pb:b2:blood-pressure:1}
\textbf{1.} $1/R = 1/7.6 + 1/22.8 + 1/4.1 + 1/5.2 + 1/5.7 + 1/19 +
1/28.5 = 0.875$: $R = \qty{1.14}{mmHg.s/mL}$; $\bar P/Q = 95/83.3 =
1.14$.
\textbf{2.} $\bar P/R_i$: encéfalo 0.75, \omterm{def:b2:heart:anatomy}{coração} 0.25, intestino e
fígado 1.39, rins 1.10, músculos 1.0, pele 0.30, outros
\qty{0.20}{L/min}; frações de 15, 5, 28, 22, 20, 6 e \qty{4}{\%}.
\textbf{3.} $5000/70 = \qty{71}{mL}$; $5\times 50 = \qty{250}{mL}$ de
oxigênio por minuto.
\textbf{4.} Encéfalo $0.75/1.4 = \qty{0.54}{L/min/kg}$; rins $1.1/0.3
= \qty{3.7}{L/min/kg}$, sete vezes o do encéfalo e cinquenta vezes a
média do corpo: os rins são irrigados para filtrar, e não para se
nutrir.
\textbf{5.} Com uma única pressão compartilhada por todos os órgãos,
cada um recebe o fluxo que sua própria resistência fixa; regular os
fluxos de modo central privaria qualquer órgão cuja demanda mudasse. A
pressão é a moeda comum.
\textbf{6.} \omterm{def:b2:heart:anatomy}{Coração}: $250\times 0.2\times 0.7 = \qty{35}{mL/min}$;
músculo: $1000\times 0.2\times 0.25 = \qty{50}{mL/min}$. O \omterm{def:b2:heart:anatomy}{coração} já
extrai a maior parte do que passa, de modo que só pode ganhar oxigênio
com mais fluxo, e é por isso que suas \omterm{def:b2:blood-circulation:vessels}{arteríolas} se dilatam no momento
em que ele trabalha mais.
\textbf{7.} $Q$ \qty{30}{\%} menor com $R$ fixo: $\bar P$ \qty{30}{\%}
menor, \qty{28.5}{mmHg}, até \qty{66.5}{mmHg}.
\textbf{8.} $28.5/5 = \qty{5.7}{mmHg}$: cerca de \qty{89}{mmHg}.
\textbf{9.} $Q = 85\times 0.7\times 71 = \qty{4.2}{L/min} =
\qty{70}{mL/s}$; $R = 89.3/70 = \qty{1.27}{mmHg.s/mL}$, um aumento de
\qty{11}{\%}.
\textbf{10.} $1060\times 9.81\times 0.4 = \qty{4160}{Pa} =
\qty{31}{mmHg}$; encéfalo $66.5 - 31 = \qty{35}{mmHg}$ antes do reflexo,
$89 - 31 = \qty{58}{mmHg}$ depois.
\textbf{11.} $28.5/2 = \qty{14}{mmHg}$: \qty{81}{mmHg} no \omterm{def:b2:heart:anatomy}{coração},
\qty{50}{mmHg} no encéfalo --- tontura, visão cinzenta, um quase
desmaio cada vez que ela se levanta de uma cadeira.
\textbf{12.} Deitar-se elimina a perda hidrostática de \qty{31}{mmHg}
e devolve ao \omterm{def:b2:heart:anatomy}{coração} o \omterm{def:b2:blood-circulation:blood}{sangue} acumulado; o \omterm{prop:b2:heart:cycle}{volume sistólico} e a pressão
no encéfalo se restabelecem em poucos batimentos.
\textbf{13.} $45/5 = \qty{9}{mmHg}$: \qty{86}{mmHg}.
\textbf{14.} $1/R = 0.875 - 1/19 - 1/4.1 + 1/38 + 1/8.2 = 0.727$: $R =
\qty{1.38}{mmHg.s/mL}$. Encéfalo $86/7.6 = \qty{11.3}{mL/s} =
\qty{0.68}{L/min}$; pele $86/38 = \qty{2.3}{mL/s} = \qty{0.14}{L/min}$;
intestino $86/8.2 = \qty{10.5}{mL/s} = \qty{0.63}{L/min}$.
\textbf{15.} Com menos \omterm{def:b2:blood-circulation:blood}{sangue} e \omterm{def:b2:blood-circulation:vessels}{arteríolas} mais contraídas, a pressão
\omterm{def:b2:blood-circulation:vessels}{capilar} cai abaixo da \omterm{thm:b2:blood-circulation:oncotic}{pressão oncótica} ao longo de todo o \omterm{def:b2:blood-circulation:vessels}{capilar}, de
modo que o líquido é reabsorvido em vez de filtrado; o \omterm{def:b2:blood-circulation:blood}{plasma} se dilui e
o \omterm{def:b2:blood-circulation:blood}{hematócrito} cai.
\textbf{16.} A aldosterona, por meio da renina e da angiotensina,
desencadeada pela baixa perfusão do rim e pela estimulação simpática; o
\omterm{prop:b2:blood-pressure:raas}{hormônio antidiurético}, desencadeado pela queda do volume e pelo aumento
da concentração plasmática.
\textbf{17.} $5\times 10^{12}$ células a $2.5\times 10^{6}$ por segundo:
$2\times 10^{6}$ s, cerca de três semanas.
\textbf{18.} Abaixo de cerca de \qty{60}{mmHg}, a curva do reflexo é
plana: os vasos estão tão contraídos e o \omterm{def:b2:heart:anatomy}{coração} tão rápido quanto
podem, o ganho é zero, e cada perda adicional baixa a pressão por
inteiro; os tecidos mal irrigados liberam então ácidos e se dilatam, e
a pressão desaba.
\textbf{19.} Músculos $110/0.33 = \qty{333}{mL/s} = \qty{20}{L/min}$;
pele $110/5 = \qty{22}{mL/s} = \qty{1.3}{L/min}$.
\textbf{20.} Resistência do encéfalo $8.8$ (fluxo de 0.75 a 110), do
\omterm{def:b2:heart:anatomy}{coração} $6.6$ (fluxo de 1.0):
$1/R = 1/8.8 + 1/6.6 + 1/8.2 + 1/10.4 + 1/0.33 + 1/5 +
1/28.5 = 3.75$: $R = \qty{0.27}{mmHg.s/mL}$; $Q = 110/0.27 =
\qty{410}{mL/s} = \qty{25}{L/min}$.
\textbf{21.} $\num{24700}/180 = \qty{137}{mL}$.
\textbf{22.} $24.7\times(200 - 40) = \qty{3950}{mL/min}$, dezesseis
vezes os \qty{250}{mL/min} de repouso.
\textbf{23.} Fluxo $\times 4.9$, extração $160/50 = \times 3.2$;
$4.9\times 3.2 = 16$.
\textbf{24.} O comando do córtex motor e os sinais dos receptores
musculares reponderam, no bulbo, a informação dos \omterm{def:b2:blood-pressure:baroreflex}{barorreceptores}, de
modo que o reflexo defende \qty{110}{mmHg} em vez de 95 --- o valor de
referência é elevado, e a alça não é desativada.
\textbf{25.} Queda residual ao se levantar de \qty{5.7}{mmHg};
\qty{86}{mmHg} após a hemorragia; \qty{25}{L/min} e \qty{4}{L} de
oxigênio por minuto no exercício.
\end{solution}
